% Part: proof-theory
% Chapter: sequent-calculus
% Section: introduction

\documentclass[../../../include/open-logic-section]{subfiles}

\begin{document}

\olfileid{pt}{seq}{int}

\olsection{Pambuka}

Kalkulus sekuen minangka sistem !!{proof} sing dikenalake dening
Gerhard Gentzen ing taun 1930-an. Ancas utamane dudu dadi cara
praktis kanggo !!{proving} samubarang. Gentzen nggunakake sistem
iki supaya bisa mbuktekake asil tartamtu ngenani struktur sing
bisa diduweni !!{proof}s lan sipat-sipate. Pitakon kaya mangkono
yaiku pitakon sing ditliti ing teori bukti.

Kalkulus sekuen tumindak marang \emph{sekuen}, manut jenenge,
dudu marang !!{formula}s. Sekuen minangka pasangan urutan utawa
multiset !!{formula}s. Sistem asli Gentzen (\Log{LK} lan \Log{LJ})
nggunakake urutan; saiki ahli teori bukti luwih milih multiset.
Multiset padha karo himpunan, nanging multiset bisa ngemot
!!{element} luwih saka sapisan. Nanging kaya ing himpunan,
urutan !!{element}s ora wigati. Dadi ekspresi $\{!A, !A, !B\}$
lan $\{!A, !B, !A\}$ nggambarake multiset sing padha, nanging
multiset kasebut beda saka $\{!A, !B\}$ lan $\{!A, !A, !B, !B\}$,
amarga loro sing pungkasan ngemot cacah $!A$ lan $!B$ sing beda
saka multiset kapisan.

Lumrahe, multiset !!{formula}s ing teori bukti ditulis nganggo
aksara Yunani gedhe. Dadi nalika kita nulis $\Gamma$, $\Delta$,
$\Pi$, $\Lambda$, utawa $\Theta$, iki dimaksudake minangka
multiset !!{formula}s ing logika sing lagi digunakake. Uga
lumrahe, ahli teori bukti ora nulis $\tuple{\Gamma, \Delta}$,
nanging misahake rong komponen nganggo tandha sekuen;
kita nggunakake~$\Sequent$.

\begin{defn}[Sekuen]
\emph{Sekuen} minangka ekspresi awujud
\[
\Gamma \Sequent \Delta
\]
kanthi $\Gamma$ lan $\Delta$ minangka multiset !!{formula}s winates
sing bisa uga kosong. $\Gamma$ diarani \emph{anteseden}, dene
$\Delta$ diarani \emph{sukseden}.
\end{defn}

Upamane $!G$~minangka konjungsi !!{formula}s ing~$\Gamma$
(konjungsi kosong dianggep~$\ltrue$), lan $!D$~minangka
disjungsi !!{formula}s ing~$\Delta$ (disjungsi kosong
dianggep~$\lfalse$). Mula sekuen $\Gamma \Sequent \Delta$ bener
(ing struktur~$\Struct{M}$) yen lan mung yen $!G \lif !D$ bener.
Ing logika klasik, iki padha karo nyatakake yen ana siji
!!{formula}s ing~$\Gamma$ sing salah, utawa ana siji
!!{formula}s ing~$\Delta$ sing bener.

Yen $\Gamma$ minangka multiset !!{formula}s, kita nulis
$\Gamma, !A$ (utawa $!A, \Gamma$) kanggo asil nambahake
$!A$ ing~$\Gamma$. Yen $\Delta$ uga multiset !!{formula}s,
maka $\Gamma, \Delta$ minangka gabungan multiset saka rong
multiset kasebut---yen $\Gamma$ ngemot $!A$ kanthi
\emph{multiplikitas}~n (yaiku ngemot $n$ salinan~$!A$), lan
$\Delta$ ngemot kanthi multiplikitas~$m$, mula
$\Gamma, \Delta$ ngemot $!A$ kanthi multiplikitas~$n+m$.
Yen multiset ing salah siji sisih sekuen kosong, lumrahe
kita mung ninggalake papan kosong. Upamane, $\Sequent !A$
yaiku sekuen kanthi anteseden kosong lan sukseden sing
ngemot $!A$ kanthi multiplikitas~$1$.

!!^a{proof} ing kalkulus sekuen minangka wit sekuen: sekuen
paling ndhuwur (utawa sekuen awal) minangka aksioma ing
sistem sing digunakake; yen sekuen ana ing sangisore siji
utawa rong sekuen liyane, sekuen kasebut kudu dijupuk kanthi
bener saka aturan inferensi sistem. Dhisik kita bakal
ngrembug rong sistem sekuen kanggo logika klasik,
$\Log{G1c}$ lan~$\Log{G3c}$. Aksioma lan aturane ana ing
\cref{pt:seq:int:tab:G1c,pt:seq:int:tab:G3c}.
\oliflabeldef{fol:seq::chap}{ (Sistem asli Gentzen \Log{LK}
diterangake ing \olref[fol][seq][]{chap}.)}{}

\subfile{rules-G1c}
\subfile{rules-G3c}

Ana prabedan alus antarane sistem iki, lan prabedan kasebut
njalari sipat teori buktine beda. \Log{G1c} duwe aksioma
$!A \Sequent !A$ tanpa !!{formula}s liyane. \Log{G3c} duwe
aksioma awujud $!C, \Gamma \Sequent \Delta, !C$, kanthi
!!{formula}s ``sisih'' sembarang ing $\Gamma$ lan~$\Delta$,
nanging !!{formula}~$!C$ sing ana ing loro sisih kudu
\emph{atomik}. \Log{G1c} duwe rong versi saben aturan
$\LeftR{\land}$ lan $\RightR{\lor}$, dene \Log{G3c} mung
duwe siji versi saben aturan. \Log{G1c} uga duwe ``aturan
struktural'' tambahan, yaiku ``kontraksi''
(\LeftR{\Contraction}, \RightR{\Contraction}) lan
``pelemahan'' (\LeftR{\Weakening}, \RightR{\Weakening}).

Sistem \Log{G1c} lan \Log{G3c} sahih lan jangkep tumrap
logika klasik. Kanthi tembung liya, saben sekuen sing
!!{provable} ing sistem iki bener ing saben !!{structure},
lan saben sekuen sing bener ing saben !!{structure} duwe
!!a{proof}. Dadi pitakon \emph{apa} sekuen duwe !!a{proof}
uga bisa diwales nganggo cara logika liyane, upamane cara
teori model. Amarga loro sistem !!{prove} persis sekuen
sing bener ing saben !!{structure}, mligine loro sistem
mbuktekake sekuen sing padha.

Nanging teori bukti ora mung ngrembug \emph{apa} samubarang
bisa dibuktekake, nanging uga \emph{kepiye} carane.
Ahli teori bukti ngrembug pitakon kaya ``Apa aturan
tartamtu mesthi bisa disingkiri?,'' ``Apa aturan iki
bisa ditambahake ing sistem tanpa ndadekake sekuen
anyar bisa dibuktekake?,'' utawa ``Apa !!{proof}s ing
siji sistem bisa diowahi dadi !!{proof}s ing sistem
liyane?'' Yen wangsulane ``iya,'' buktine asring
nggunakake induksi tumrap kerumitan !!{proof}s,
utawa, kanthi ekuivalen, induksi tumrap struktur
!!{proof}s. Asile ora mung bukti tumrap pratelan
kasebut (upamane, yen \Log{G1c} mbuktekake sekuen,
maka \Log{G3c} mbuktekake sekuen sing padha), nanging
uga prosedur (upamane, ngowahi !!{proof}s ing~\Log{G1c}
dadi !!{proof}s ing~\Log{G3c}).

Asil utama ing teori bukti yaiku \emph{teorema eliminasi
potongan}. Teorema iki nuduhake yen kita bisa nambahake
aturan potongan
\[
\Axiom$\Gamma \fCenter \Delta, !A$
\Axiom$!A, \Pi \fCenter \Lambda$
\RightLabel{\Cut}
\BinaryInf$\Gamma, \Pi \fCenter \Delta, \Lambda$
\DisplayProof
\]
ing sistem sekuen tanpa ngowahi sekuen endi sing
!!{provable}. Kajaba iku, buktine menehi cara ngowahi
saben !!{proof} sing nggunakake aturan \Log{G1c}
(utawa \Log{G3c}) bebarengan karo \Cut{} dadi bukti
sing mung nggunakake aturan \Log{G1c} (utawa \Log{G3c}).
Kanthi tembung liya, prosedur iki ``ngilangi'' aturan
\Cut{} saka !!{proof}s sing nggunakake aturan kasebut;
mula teorema iki diwenehi jeneng mangkono.

\end{document}
